\documentclass[11pt,a4paper]{article}

\usepackage[utf8]{inputenc}
\usepackage[T1]{fontenc}
\usepackage[spanish,es-noquoting]{babel}
\usepackage{lmodern}
\usepackage[margin=2.3cm]{geometry}
\usepackage{xcolor}
\usepackage{listings}
\usepackage{booktabs}
\usepackage{array}
\usepackage{enumitem}
\usepackage{tikz}
\usetikzlibrary{arrows.meta,positioning,automata,shapes.geometric,calc,decorations.pathreplacing,fit,backgrounds}
\usepackage[most]{tcolorbox}
\usepackage{fancyhdr}
\usepackage[hidelinks]{hyperref}

% ---------------- Paleta ----------------
\definecolor{tinta}{RGB}{30,41,59}
\definecolor{acento}{RGB}{0,110,130}
\definecolor{cBoton}{RGB}{230,126,0}     % botón
\definecolor{cPulso}{RGB}{200,40,70}     % pulso
\definecolor{cCaida}{RGB}{34,140,80}     % caída / acelerómetro
\definecolor{cEsp}{RGB}{60,80,160}       % ESP32
\definecolor{cPC}{RGB}{110,70,160}       % computador
\definecolor{cLed}{RGB}{220,170,0}       % LED
\definecolor{cGnd}{RGB}{40,40,40}
\definecolor{c3v3}{RGB}{215,38,61}
\definecolor{fondo}{RGB}{247,248,250}
\definecolor{clave}{RGB}{0,82,155}
\definecolor{coment}{RGB}{100,120,100}
\definecolor{cadena}{RGB}{163,21,21}

\lstset{
  basicstyle=\ttfamily\footnotesize,
  backgroundcolor=\color{fondo},
  keywordstyle=\color{clave}\bfseries,
  commentstyle=\color{coment}\itshape,
  stringstyle=\color{cadena},
  numbers=left, numberstyle=\tiny\color{gray}, numbersep=8pt,
  frame=single, rulecolor=\color{gray!40},
  breaklines=true, showstringspaces=false, tabsize=2,
  xleftmargin=1.5em, framexleftmargin=1.2em,
  literate={á}{{\'a}}1 {é}{{\'e}}1 {í}{{\'i}}1 {ó}{{\'o}}1 {ú}{{\'u}}1
           {Á}{{\'A}}1 {É}{{\'E}}1 {Í}{{\'I}}1 {Ó}{{\'O}}1 {Ú}{{\'U}}1
           {ñ}{{\~n}}1 {¡}{{!`}}1 {¿}{{?`}}1 {Ω}{{$\Omega$}}1 {±}{{$\pm$}}1
}
\lstdefinelanguage{ini}{morecomment=[l]{;}, morekeywords={platform,board,framework,monitor_speed,build_flags,monitor_rts,monitor_dtr}}

\newtcolorbox{nota}[1][]{colback=acento!6,colframe=acento,fonttitle=\bfseries,title=#1,arc=2pt,boxrule=0.6pt,before upper={\setlength{\parskip}{4pt}}}
\newtcolorbox{analogia}[1][]{colback=cEsp!6,colframe=cEsp,fonttitle=\bfseries,title={Analogía con software: #1},arc=2pt,boxrule=0.6pt,before upper={\setlength{\parskip}{4pt}}}
\newtcolorbox{ojo}[1][]{colback=cPulso!6,colframe=cPulso,fonttitle=\bfseries,title=#1,arc=2pt,boxrule=0.6pt}

\newcommand{\chip}[2]{{\setlength{\fboxsep}{1.6pt}\fcolorbox{#1}{#1!15}{\textcolor{#1!70!black}{\scriptsize\bfseries\ttfamily #2}}}}

\pagestyle{fancy}
\fancyhf{}
\lhead{\small Alerta ESP32-S3}
\rhead{\small Funcionamiento del sistema}
\cfoot{\thepage}
\setlength{\headheight}{14pt}

\title{\textbf{Alerta ESP32-S3}\\[6pt]\large Funcionamiento del sistema: hardware, lógica y comunicación}
\author{Sara Otero}
\date{\today}

\begin{document}
\shorthandoff{"}
\maketitle

\begin{nota}[A quién va dirigido]
Este documento explica cómo funciona el dispositivo de alerta para personas con
formación en sistemas. La parte de software (estados, HTTP, JSON, colas) se presenta
de forma directa; la parte electrónica (voltajes, sensores, buses de comunicación)
se explica desde cero, con analogías a conceptos de software cuando ayudan.
\end{nota}

\tableofcontents
\newpage

%=====================================================================
\section{Visión general}
%=====================================================================

El dispositivo es un \textbf{sistema embebido} (un computador pequeño dedicado a una
sola tarea) que vigila tres posibles situaciones de emergencia de la persona que lo
lleva y avisa a un computador por WiFi:

\begin{center}
\begin{tabular}{@{}lll@{}}
\toprule
\textbf{Causa} & \textbf{Qué la dispara} & \textbf{Sensor}\\
\midrule
\chip{cBoton}{BOTON} & La persona oprime el botón de pánico & Pulsador\\
\chip{cPulso}{RITMO} & El ritmo cardíaco llega a 110 pulsaciones por minuto o más & XD-58C\\
\chip{cCaida}{CAIDA} & Se detecta una caída (caída libre seguida de un golpe) & MPU6050\\
\bottomrule
\end{tabular}
\end{center}

Cualquiera de las tres abre una \textbf{ventana de 30 segundos} en la que la persona
puede cancelar con el botón (por si fue un falso positivo). Si nadie cancela, se envía
la alerta al computador. El dispositivo nunca deja de vigilar: después de enviar una
alerta sigue leyendo los sensores, y cada evento nuevo genera su propia alerta.

\subsection{Arquitectura}

\begin{center}
\begin{tikzpicture}[
  font=\small,
  sen/.style={draw,rounded corners=4pt,minimum width=2.9cm,minimum height=1cm,align=center,line width=0.8pt},
  flecha/.style={-{Stealth[length=3mm]},line width=1pt}]
  \node[sen,draw=cBoton,fill=cBoton!12] (b) at (0,2.2) {Botón\\\scriptsize entrada digital};
  \node[sen,draw=cPulso,fill=cPulso!12] (p) at (0,0.7) {Sensor de pulso\\\scriptsize XD-58C, analógico};
  \node[sen,draw=cCaida,fill=cCaida!12] (a) at (0,-0.8) {Acelerómetro\\\scriptsize MPU6050, bus I2C};
  \node[sen,draw=cLed,fill=cLed!15] (l) at (0,-2.3) {LED\\\scriptsize salida digital};
  \node[draw=cEsp,fill=cEsp!10,rounded corners=6pt,line width=1pt,minimum width=3.6cm,minimum height=5.4cm,align=center] (e) at (4.9,0) {\textbf{ESP32-S3}\\[2pt]\scriptsize microcontrolador\\[8pt]\scriptsize \texttt{src/main.cpp}\\\scriptsize máquina de estados\\\scriptsize + lectura de sensores\\\scriptsize + cola de envíos};
  \node[draw=cPC,fill=cPC!10,rounded corners=6pt,line width=1pt,minimum width=3.2cm,minimum height=2.2cm,align=center] (pc) at (11.6,0) {\textbf{Computador}\\[2pt]\scriptsize \texttt{pc/servidor\_pc.py}\\\scriptsize servidor HTTP};
  \draw[flecha,cBoton] (b.east) -- (b.east -| e.west);
  \draw[flecha,cPulso] (p.east) -- (p.east -| e.west);
  \draw[flecha,cCaida,{Stealth[length=3mm]}-{Stealth[length=3mm]}] (a.east) -- (a.east -| e.west);
  \draw[flecha,cLed] (l.east -| e.west) -- (l.east);
  \draw[flecha,cPC,dashed] (e.east) -- node[above]{\scriptsize WiFi} node[below]{\scriptsize HTTP POST + JSON} (pc.west);
\end{tikzpicture}
\end{center}

Hay dos programas que corren en lugares distintos:

\begin{itemize}[leftmargin=*]
  \item \textbf{Firmware} (el programa que vive grabado en el microcontrolador):
        \texttt{src/main.cpp}, en C++ con el framework de Arduino. Lee los sensores,
        decide y envía.
  \item \textbf{Servidor} en el computador: \texttt{pc/servidor\_pc.py}, en Python sin
        dependencias externas. Recibe las alertas y las muestra.
\end{itemize}

%=====================================================================
\section{Conceptos de electrónica necesarios}
%=====================================================================

Esta sección resume lo mínimo de electrónica para entender el resto del documento.

\subsection{Voltaje, 3V3 y GND}

El \textbf{voltaje} es la ``presión'' eléctrica entre dos puntos, medida en voltios (V).
En este circuito hay dos referencias que aparecen en todas partes:

\begin{itemize}[leftmargin=*]
  \item \textbf{3V3} (3.3 V): la alimentación. Es lo que la ESP32 le entrega a los sensores
        para que funcionen.
  \item \textbf{GND} (\emph{ground}, tierra): la referencia de 0 V. Es el camino de regreso
        de la corriente. Todos los componentes deben compartir el mismo GND; si no, sus
        voltajes no son comparables.
\end{itemize}

\begin{analogia}[GND común]
Es como un sistema de coordenadas compartido: un valor de ``3.3 V'' solo tiene sentido
si todos miden desde el mismo cero. Dos dispositivos sin GND común son como dos
servicios con relojes no sincronizados comparando \emph{timestamps}.
\end{analogia}

\subsection{GPIO: los pines del microcontrolador}

Un \textbf{GPIO} (\emph{General Purpose Input/Output}) es un pin que el programa puede
usar como entrada o como salida. Se identifican por número: GPIO4, GPIO2, etc.

\begin{center}
\begin{tabular}{@{}>{\bfseries}l p{11cm}@{}}
\toprule
Modo & Qué significa\\
\midrule
Salida digital & El programa escribe \texttt{HIGH} (3.3 V) o \texttt{LOW} (0 V). Así se prende y apaga el LED.\\
Entrada digital & El programa lee si el pin está en \texttt{HIGH} o \texttt{LOW}. Así se lee el botón.\\
Entrada analógica & El programa lee un \emph{voltaje continuo} y lo convierte en un número (ver ADC). Así se lee el sensor de pulso.\\
Bus de datos (I2C) & Dos pines que juntos forman un canal de comunicación con otro chip. Así se habla con el acelerómetro.\\
\bottomrule
\end{tabular}
\end{center}

\subsection{ADC: convertir voltaje en números}

El \textbf{ADC} (\emph{Analog to Digital Converter}) mide un voltaje y lo traduce a un entero.
En la ESP32-S3 la resolución es de 12 bits:

\begin{center}
\begin{tikzpicture}[font=\small]
  \shade[left color=cEsp!5,right color=cEsp!45] (0,0) rectangle (10,0.6);
  \draw (0,0) rectangle (10,0.6);
  \foreach \x/\v/\n in {0/0 V/0, 5/1.65 V/2048, 10/3.3 V/4095} {
    \draw (\x,0) -- (\x,-0.15);
    \node[below] at (\x,-0.15) {\texttt{\n}};
    \node[above] at (\x,0.6) {\v};
  }
  \node[left] at (-0.2,0.3) {voltaje};
  \node[left] at (-0.2,-0.45) {número leído};
\end{tikzpicture}
\end{center}

La función \texttt{analogRead(PIN\_PULSO)} devuelve ese número, entre 0 y 4095.

\subsection{Pull-up: por qué el botón ``presionado'' es LOW}

Un pin de entrada que no está conectado a nada queda \textbf{flotando}: lee valores
aleatorios porque capta ruido eléctrico. Para evitarlo se usa una resistencia de
\textbf{pull-up}, que conecta el pin débilmente a 3.3 V. La ESP32 trae una interna, que
se activa con \texttt{pinMode(PIN\_BOTON, INPUT\_PULLUP)}.

\begin{center}
\begin{tikzpicture}[font=\small,thick]
  \node[draw=cEsp,fill=cEsp!10,rounded corners,minimum width=2.4cm,minimum height=1.8cm] (esp) at (0,0) {GPIO4};
  \draw (1.2,0) -- (3,0) coordinate (n);
  \draw (n) -- ++(0,0.5);
  \draw[c3v3] (3,0.5) -- (3,0.7) -- (2.85,0.8) -- (3.15,0.95) -- (2.85,1.1) -- (3.15,1.25) -- (3,1.35) -- (3,1.6);
  \node[c3v3,right] at (3.15,1.05) {\scriptsize pull-up interno};
  \node[c3v3,above] at (3,1.6) {3V3};
  \draw (n) -- (4.5,0);
  \draw (4.5,0) -- (4.5,-0.2); \draw (4.3,-0.2) -- (4.75,-0.55); \draw (4.5,-0.65) -- (4.5,-1);
  \node[right] at (4.8,-0.4) {\scriptsize botón};
  \draw (4.3,-1) -- (4.7,-1); \draw (4.37,-1.1) -- (4.63,-1.1); \draw (4.44,-1.2) -- (4.56,-1.2);
  \node[below] at (4.5,-1.2) {\scriptsize GND};
  \node[align=left,anchor=west,font=\small] at (6.4,0.5) {Suelto: el pull-up lleva el pin a 3.3 V \(\rightarrow\) \texttt{HIGH}};
  \node[align=left,anchor=west,font=\small] at (6.4,-0.5) {Presionado: el botón lo une a GND \(\rightarrow\) \texttt{LOW}};
\end{tikzpicture}
\end{center}

Por eso en el código ``presionado'' es \texttt{LOW}: la lógica está invertida respecto a
lo intuitivo.

\subsection{Rebote mecánico (\emph{debounce})}

Un botón físico no cambia de estado limpiamente: durante unos milisegundos sus contactos
metálicos rebotan y el pin alterna rápido entre \texttt{HIGH} y \texttt{LOW}.

\begin{center}
\begin{tikzpicture}[font=\small,thick,xscale=1.1]
  \draw[->,gray] (0,0) -- (10,0) node[right,black] {tiempo};
  \draw[cBoton,line width=1.2pt] (0,1) -- (2,1) -- (2,0) -- (2.2,0) -- (2.2,1) -- (2.35,1) -- (2.35,0) -- (2.6,0) -- (2.6,1) -- (2.7,1) -- (2.7,0) -- (6,0) -- (6,1) -- (6.15,1) -- (6.15,0) -- (6.3,0) -- (6.3,1) -- (10,1);
  \node[left] at (0,1) {\texttt{HIGH}};
  \node[left] at (0,0) {\texttt{LOW}};
  \draw[decorate,decoration={brace,amplitude=4pt}] (2,1.2) -- node[above=4pt]{\scriptsize rebotes ($<$ 50 ms)} (2.75,1.2);
  \node[cBoton] at (4.4,0.5) {\scriptsize presionado};
\end{tikzpicture}
\end{center}

Sin tratamiento, una sola pulsación se leería como varias (y activaría y cancelaría la
alerta en el mismo instante). La solución, llamada \textbf{debounce}, es aceptar un cambio
solo cuando la lectura lleva más de 50 ms estable.

\begin{analogia}[debounce]
Es el mismo concepto que el \emph{debounce} de eventos en una interfaz web: ignorar
ráfagas de eventos y reaccionar solo cuando la señal se estabiliza.
\end{analogia}

\subsection{I2C: el bus de comunicación del acelerómetro}

El acelerómetro no entrega un voltaje simple: es un chip con su propio procesador y
registros internos. Se comunica por \textbf{I2C}, un bus serial de dos cables:

\begin{itemize}[leftmargin=*]
  \item \textbf{SDA} (\emph{Serial Data}): por donde viajan los datos, bit a bit.
  \item \textbf{SCL} (\emph{Serial Clock}): una señal de reloj que marca cuándo leer cada bit.
\end{itemize}

Cada chip en el bus tiene una dirección; la del MPU6050 es \texttt{0x68}. Leer la
aceleración consiste en pedirle al chip el contenido de unos registros (posiciones de
memoria) a partir de la dirección \texttt{0x3B}.

\begin{analogia}[I2C]
Es parecido a un protocolo petición-respuesta sobre un bus compartido: la ESP32 actúa
como cliente, el sensor como servidor con dirección fija, y los ``registros'' son como
\emph{endpoints} de lectura y escritura. Si un cable está flojo, el resultado es el
equivalente a un \emph{timeout} o una respuesta corrupta.
\end{analogia}

\subsection{Aceleración en ``g''}

El acelerómetro mide aceleración en los ejes $x$, $y$, $z$. La unidad práctica es la
\textbf{g}, la aceleración de la gravedad (9.8 m/s$^2$). El programa calcula la magnitud
total:
\[
  |a| = \sqrt{a_x^2 + a_y^2 + a_z^2}
\]

\begin{center}
\begin{tabular}{@{}cl@{}}
\toprule
\textbf{Valor} & \textbf{Situación}\\
\midrule
$\approx 1$ g & Quieto (solo se siente la gravedad), sin importar la orientación\\
$\approx 0$ g & Cayendo libremente (en caída libre no se siente peso)\\
$> 2$ g & Golpe fuerte (el cuerpo se frena de golpe al tocar el suelo)\\
\bottomrule
\end{tabular}
\end{center}

%=====================================================================
\section{Hardware y conexiones}
%=====================================================================

\subsection{Componentes}

\begin{center}
\begin{tabular}{@{}l p{10.2cm}@{}}
\toprule
\textbf{Componente} & \textbf{Función}\\
\midrule
ESP32-S3 DevKitC-1 & Microcontrolador con WiFi. Ejecuta el firmware.\\
Pulsador (botón) & Botón de pánico y botón para cancelar.\\
LED + resistencia 220 $\Omega$ & Indicador de estado. La resistencia limita la corriente para no quemar el LED.\\
XD-58C & Sensor óptico de pulso: un LED verde ilumina la piel y un fotosensor mide cuánta luz vuelve; cambia con cada latido. Entrega un voltaje analógico.\\
MPU6050 (módulo GY-521) & Acelerómetro y giroscopio de 3 ejes. Se usa solo el acelerómetro.\\
\bottomrule
\end{tabular}
\end{center}

\subsection{Tabla de conexiones}

\begin{center}
\begin{tabular}{@{}llll@{}}
\toprule
\textbf{Componente} & \textbf{Pin del componente} & \textbf{Pin de la ESP32} & \textbf{Tipo de señal}\\
\midrule
\color{cBoton}Botón & una pata & GPIO4 & entrada digital (pull-up)\\
 & pata opuesta (diagonal) & GND & \\
\color{cLed!80!black}LED & ánodo (pata larga) vía 220 $\Omega$ & GPIO2 & salida digital\\
 & cátodo (pata corta) & GND & \\
\color{cPulso}XD-58C & S (señal) & GPIO5 & entrada analógica (ADC)\\
 & + & 3V3 & alimentación\\
 & $-$ & GND & \\
\color{cCaida}MPU6050 & SDA & GPIO8 & bus I2C (datos)\\
 & SCL & GPIO9 & bus I2C (reloj)\\
 & VCC & 3V3 & alimentación\\
 & GND & GND & \\
\bottomrule
\end{tabular}
\end{center}

\begin{center}
\begin{tikzpicture}[font=\small,thick]
  \node[draw=cEsp,fill=cEsp!10,rounded corners=6pt,minimum width=3.2cm,minimum height=6cm,line width=1pt] (esp) at (0,0) {};
  \node[cEsp!80!black,font=\bfseries] at (0,2.6) {ESP32-S3};
  \foreach \y/\t in {2/GPIO4, 1.2/GPIO2, 0.4/GPIO5, -0.4/GPIO8, -1.2/GPIO9, -2/3V3, -2.6/GND} {
    \node[anchor=east,font=\scriptsize\ttfamily] at (1.55,\y) {\t};
    \fill (1.6,\y) circle (2pt);
  }
  % Botón
  \draw[cBoton,line width=1.4pt] (1.6,2) -- (4,2);
  \node[draw=cBoton,fill=cBoton!12,rounded corners,minimum width=2.6cm] at (5.4,2) {Botón $\rightarrow$ GND};
  % LED
  \draw[cLed!80!black,line width=1.4pt] (1.6,1.2) -- (4,1.2);
  \node[draw=cLed,fill=cLed!15,rounded corners,minimum width=2.6cm] at (5.4,1.2) {220 $\Omega$ + LED $\rightarrow$ GND};
  % Pulso
  \draw[cPulso,line width=1.4pt] (1.6,0.4) -- (4,0.4);
  \node[draw=cPulso,fill=cPulso!12,rounded corners,minimum width=2.6cm] at (5.4,0.4) {XD-58C (S)};
  % MPU
  \draw[cCaida,line width=1.4pt] (1.6,-0.4) -- (4,-0.4);
  \draw[cCaida,line width=1.4pt,dashed] (1.6,-1.2) -- (4,-1.2);
  \node[draw=cCaida,fill=cCaida!12,rounded corners,minimum width=2.6cm,minimum height=1.2cm,align=center] at (5.4,-0.8) {MPU6050\\\scriptsize SDA (—) y SCL (- -)};
  % Alimentación
  \draw[c3v3,line width=1.4pt] (1.6,-2) -- (8,-2) node[right] {riel + (3V3): XD-58C y MPU6050};
  \draw[cGnd,line width=1.4pt] (1.6,-2.6) -- (8,-2.6) node[right] {riel $-$ (GND): todos};
\end{tikzpicture}
\end{center}

\begin{ojo}[Cuidado con el voltaje]
Los sensores se alimentan con 3V3, nunca con 5 V. Los pines de la ESP32 toleran como
máximo 3.3 V; si el sensor de pulso se alimenta con 5 V, su salida puede superar ese
límite y dañar el pin.
\end{ojo}

\begin{ojo}[Conexiones flojas]
El bus I2C y la señal analógica son sensibles a contactos intermitentes. Con pines sin
soldar aparecen errores de lectura del acelerómetro y ruido en el pulso. El firmware
tolera estas fallas (sección~\ref{sec:robustez}), pero la solución definitiva es soldar
los pines.
\end{ojo}

%=====================================================================
\section{Lógica del sistema}
%=====================================================================

\subsection{Máquina de estados}

El firmware tiene dos estados. La lectura de los sensores y el envío de alertas ocurren
\textbf{en ambos estados}, en cada iteración del ciclo principal.

\begin{center}
\begin{tikzpicture}[->,>={Stealth[length=3mm]},node distance=6.5cm,thick,
  every state/.style={minimum size=2.6cm,align=center,line width=1.2pt}]
  \node[state,initial,initial text={arranque},draw=cEsp,fill=cEsp!10] (E) {\textbf{ESTABLE}\\\scriptsize vigilando\\\scriptsize LED titila};
  \node[state,right of=E,draw=cPulso,fill=cPulso!10] (P) {\textbf{PENDIENTE}\\\scriptsize contando 30 s\\\scriptsize LED fijo};
  \path (E) edge[bend left=28] node[above=3pt,align=center,font=\small] {\chip{cBoton}{BOTON} o \chip{cPulso}{RITMO} o \chip{cCaida}{CAIDA}} (P)
        (P) edge[bend left=28] node[below,font=\small] {botón (cancelar)} (E)
        (P) edge[bend left=55,looseness=1] node[below,font=\small,align=center] {pasan 30 s:\\ alerta a la cola de envío} (E)
        (P) edge[loop right] node[right,font=\small,align=center] {otra causa:\\se suma} (P);
\end{tikzpicture}
\end{center}

\begin{center}
\begin{tabular}{@{}l p{10.6cm}@{}}
\toprule
\textbf{Transición} & \textbf{Comportamiento}\\
\midrule
ESTABLE $\rightarrow$ PENDIENTE & Ocurre una o varias causas en la misma iteración. Se guardan como un conjunto de \emph{flags} y empieza el conteo de 30 s.\\
PENDIENTE $\rightarrow$ PENDIENTE & Ocurre otra causa durante el conteo: se agrega al conjunto (por ejemplo \texttt{RITMO+CAIDA}). El conteo no se reinicia.\\
PENDIENTE $\rightarrow$ ESTABLE & El botón cancela: no se envía nada.\\
PENDIENTE $\rightarrow$ ESTABLE & Se cumplen los 30 s: la alerta pasa a la cola de envío y el sistema vuelve a vigilar inmediatamente.\\
\bottomrule
\end{tabular}
\end{center}

\subsection{Indicador LED}

El LED comunica el estado sin necesidad de un computador:

\begin{center}
\begin{tikzpicture}[font=\small,x=0.42cm,y=0.5cm]
  % fila de tiempo: 1 celda = 0.5 s
  \def\onC{cLed}
  \foreach \i in {0,...,7} { \pgfmathparse{mod(\i,2)==0 ? 1 : 0}\ifnum\pgfmathresult=1 \fill[cLed] (\i,0) rectangle ++(1,1); \else \fill[gray!20] (\i,0) rectangle ++(1,1); \fi }
  \fill[cLed] (8,0) rectangle (20,1);
  \fill[gray!20] (20,0) rectangle (22,1);
  \foreach \i in {22,...,29} { \pgfmathparse{mod(\i,2)==0 ? 1 : 0}\ifnum\pgfmathresult=1 \fill[cLed] (\i,0) rectangle ++(1,1); \else \fill[gray!20] (\i,0) rectangle ++(1,1); \fi }
  \draw (0,0) grid (30,1);
  \draw[decorate,decoration={brace,amplitude=4pt}] (0,1.2) -- node[above=4pt]{vigilando: titila} (8,1.2);
  \draw[decorate,decoration={brace,amplitude=4pt}] (8,1.2) -- node[above=4pt]{contando 30 s: fijo} (20,1.2);
  \draw[decorate,decoration={brace,amplitude=4pt}] (22,1.2) -- node[above=4pt]{vigilando de nuevo: titila} (30,1.2);
  \draw[->,cPulso,thick] (8,-0.9) node[below,font=\scriptsize]{detecta algo} -- (8,-0.1);
  \draw[->,cPC,thick] (21,-0.9) node[below,font=\scriptsize,align=center]{se envía:\\LED apagado 1 s} -- (21,-0.1);
\end{tikzpicture}
\end{center}

\begin{center}
\begin{tabular}{@{}l p{9.4cm}@{}}
\toprule
\textbf{LED} & \textbf{Significado}\\
\midrule
Titila (cada 0.5 s) & Vigilando, sin eventos\\
Fijo & Contando los 30 s; se puede cancelar con el botón\\
Apagado 1 s & Se acaba de enviar una alerta\\
Parpadeo durante el arranque & Conectándose al WiFi (fijo 2 s = conectado; parpadeo rápido = falló)\\
\bottomrule
\end{tabular}
\end{center}

\subsection{Ciclo de vida de una alerta}

\begin{center}
\begin{tikzpicture}[font=\small,node distance=0.55cm,
  paso/.style={draw,rounded corners=4pt,minimum width=9.5cm,align=center,inner sep=6pt,line width=0.8pt},
  flecha/.style={-{Stealth[length=2.5mm]},thick}]
  \node[paso,draw=cCaida,fill=cCaida!10] (s1) {Un sensor detecta un evento (por ejemplo, una caída)};
  \node[paso,draw=cPulso,fill=cPulso!10,below=of s1] (s2) {Estado PENDIENTE: LED fijo, empieza el conteo de 30 s};
  \node[paso,draw=cPulso,fill=cPulso!5,below=of s2] (s3) {Durante el conteo, otros eventos se suman a la misma alerta};
  \node[paso,draw=cBoton,fill=cBoton!10,below left=0.6cm and -3.2cm of s3,minimum width=4.4cm] (c) {Botón: se cancela\\\scriptsize vuelve a vigilar, no se envía};
  \node[paso,draw=cEsp,fill=cEsp!10,below right=0.6cm and -3.2cm of s3,minimum width=4.4cm] (q) {30 s sin cancelar:\\\scriptsize la alerta entra a la cola};
  \node[paso,draw=cPC,fill=cPC!10,below=2.7cm of s3] (s5) {Envío en segundo plano: HTTP POST al computador\\\scriptsize reintenta cada 3 s, máximo 5 intentos};
  \node[paso,draw=cEsp,fill=cEsp!5,below=of s5] (s6) {Mientras tanto el sistema ya volvió a ESTABLE y sigue vigilando};
  \draw[flecha] (s1) -- (s2);
  \draw[flecha] (s2) -- (s3);
  \draw[flecha] (s3.south) -- (c.north);
  \draw[flecha] (s3.south) -- (q.north);
  \draw[flecha] (q.south) -- (q.south |- s5.north);
  \draw[flecha] (s5) -- (s6);
\end{tikzpicture}
\end{center}

\paragraph{Ejemplo.} El pulso sube a 120 BPM: se abre un conteo por \chip{cPulso}{RITMO}.
A los 10 s la persona se cae: la alerta pasa a ser \texttt{RITMO+CAIDA}. Al cumplirse
los 30 s se envía esa alerta combinada. Si minutos después ocurre otra caída, se abre un
conteo nuevo y se envía otra alerta, esta vez solo \chip{cCaida}{CAIDA}.

%=====================================================================
\section{Los sensores en detalle}
%=====================================================================

\subsection{Sensor de pulso XD-58C}

\paragraph{Principio físico.} El sensor ilumina la piel con un LED verde y mide la luz
reflejada. Con cada latido entra más sangre a los capilares, que absorbe más luz. El
resultado es un voltaje que sube y baja al ritmo del corazón, alrededor de un valor
central de unos 1.65 V (lectura de ADC $\approx$ 2048).

\begin{center}
\begin{tikzpicture}[font=\small,x=1.2cm,y=1.6cm]
  \draw[->,gray] (0,0) -- (10.3,0) node[right,black] {tiempo};
  \draw[->,gray] (0,0) -- (0,2.3) node[above,black] {lectura ADC};
  % onda
  \draw[cPulso,line width=1.3pt,smooth] plot coordinates {
    (0,0.9) (0.3,0.9) (0.5,1.9) (0.7,1.1) (0.9,1.25) (1.1,0.95) (1.5,0.9) (2.3,0.9)
    (2.5,1.95) (2.7,1.1) (2.9,1.25) (3.1,0.95) (3.5,0.9) (4.3,0.9)
    (4.5,1.9) (4.7,1.1) (4.9,1.25) (5.1,0.95) (5.5,0.9) (6.3,0.9)
    (6.5,1.95) (6.7,1.1) (6.9,1.25) (7.1,0.95) (7.5,0.9) (8.3,0.9)
    (8.5,1.9) (8.7,1.1) (8.9,1.25) (9.1,0.95) (9.5,0.9) (10,0.9)};
  \draw[dashed,cEsp] (0,1.05) -- (10,1.05) node[right,font=\scriptsize] {línea base};
  \draw[dashed,cCaida] (0,1.45) -- (10,1.45) node[right,font=\scriptsize] {umbral de latido};
  \foreach \x in {0.42,2.42,4.42,6.42,8.42} \fill[cCaida] (\x,1.45) circle (2.2pt);
  \draw[<->,thick] (2.42,2.15) -- node[above]{\scriptsize IBI (ms entre latidos)} (4.42,2.15);
\end{tikzpicture}
\end{center}

\paragraph{Algoritmo.} Cada 10 ms se toma una muestra y se actualizan dos promedios
móviles exponenciales:
\begin{itemize}[leftmargin=*]
  \item \textbf{línea base}: el valor promedio de la señal (se adapta lentamente);
  \item \textbf{amplitud}: el tamaño típico de la oscilación alrededor de la base.
\end{itemize}
Un \textbf{latido} se registra cuando la señal cruza hacia arriba el umbral
$\text{base} + 0.5 \cdot \text{amplitud}$, con un período refractario de 300 ms (no se
aceptan dos latidos más juntos que eso). El tiempo entre dos latidos se llama
\textbf{IBI} (\emph{inter-beat interval}). Con los últimos cinco IBI:
\[
  \text{BPM} = \frac{60\,000 \cdot 5}{\sum_{i=1}^{5} \text{IBI}_i}
\]

\paragraph{Disparo.} El ritmo dispara cuando el BPM llega a \texttt{UMBRAL\_BPM}
(110). Para no disparar en bucle mientras el pulso sigue alto, el disparo se
\emph{desarma} y solo se \emph{rearma} cuando el pulso baja de 105 BPM (histéresis).

\begin{analogia}[histéresis]
Es como un \emph{rate limiter} con condición de reinicio: después de emitir el evento
no se vuelve a emitir hasta que la condición se haya apagado claramente, lo que evita
oscilaciones cuando el valor ronda el umbral.
\end{analogia}

\paragraph{Filtros anti-ruido (configurables).} Si el sensor está mal conectado, el pin
analógico queda flotando y el ruido eléctrico puede parecer latidos. El firmware tiene
filtros que se activan o desactivan desde \texttt{config.h}. En la configuración actual
están \textbf{desactivados}, para que el sensor siempre entregue una lectura durante
las pruebas con pines sin soldar.

\begin{center}
\begin{tabular}{@{}l p{5.6cm} c c@{}}
\toprule
\textbf{Filtro} & \textbf{Qué hace} & \textbf{Actual} & \textbf{Recomendado}\\
\midrule
\texttt{PULSO\_BASE\_MIN/MAX} & Descarta la señal si el promedio está pegado arriba o abajo (sensor sin alimentación) & 0 / 4096 & 300 / 3800\\
\texttt{PULSO\_REGULARIDAD} & Descarta si el IBI más largo supera en más de este factor al más corto (latidos disparejos = ruido) & 100 & 1.3\\
\texttt{IBI\_MIN/MAX\_MS} & Rango fisiológico aceptado & 300 / 2000 & 333 / 1500\\
\texttt{RITMO\_CONFIRMA\_MS} & Tiempo que el pulso debe sostenerse alto antes de disparar & 0 & 5000\\
\bottomrule
\end{tabular}
\end{center}

\subsection{Acelerómetro MPU6050: detección de caídas}

Una caída tiene una firma característica en la magnitud de la aceleración: primero un
tramo de \textbf{caída libre} (la lectura baja hacia 0 g) y, poco después, un
\textbf{impacto} (un pico por encima de 2 g).

\begin{center}
\begin{tikzpicture}[font=\small,x=1.15cm,y=1.1cm]
  \fill[cCaida!10] (2.9,0) rectangle (5.3,3.6);
  \node[cCaida!70!black,font=\scriptsize,below] at (4.1,0) {ventana de 1 s};
  \draw[->,gray] (0,0) -- (10.3,0) node[right,black] {tiempo};
  \draw[->,gray] (0,0) -- (0,3.7) node[above,black] {$|a|$ (g)};
  \foreach \y in {1,2,3} \node[left,font=\scriptsize] at (0,\y) {\y};
  \draw[dashed,cEsp] (0,0.6) -- (10,0.6) node[right,font=\scriptsize] {0.6 g: caída libre};
  \draw[dashed,cPulso] (0,2) -- (10,2) node[right,font=\scriptsize] {2.0 g: impacto};
  \draw[cCaida,line width=1.3pt] plot[smooth] coordinates {(0,1) (2.5,1.02) (2.8,0.9) (3,0.3) (3.6,0.15) (4.2,0.2) (4.4,3.2) (4.6,0.6) (4.9,1.4) (5.3,0.9) (6,1.05) (10,1)};
  \fill[cEsp] (3,0.3) circle (2.5pt);
  \fill[cPulso] (4.4,3.2) circle (2.5pt);
  \node[right,font=\scriptsize,align=left] at (5.4,3.0) {$\leftarrow$ CAIDA detectada};
  \node[below,font=\scriptsize] at (1.3,0.95) {quieto (1 g)};
\end{tikzpicture}
\end{center}

\paragraph{Algoritmo.} Cada 10 ms se lee la aceleración por I2C y se calcula $|a|$.
\begin{enumerate}[leftmargin=*]
  \item Si $|a| <$ \texttt{CAIDA\_LIBRE\_G} (0.6 g), se marca el inicio de una posible caída.
  \item Si dentro de los siguientes \texttt{VENTANA\_IMPACTO\_MS} (1 s) aparece
        $|a| >$ \texttt{IMPACTO\_G} (2.0 g), se reporta una caída.
  \item Si pasa 1 s sin impacto, se descarta (era solo un movimiento brusco).
\end{enumerate}

El sensor se configura en el rango de $\pm 8$ g, en el que cada g equivale a 4096
unidades del valor crudo de 16 bits.

\subsection{Robustez ante conexiones defectuosas}\label{sec:robustez}

\begin{center}
\begin{tabular}{@{}p{4.6cm} p{11cm}@{}}
\toprule
\textbf{Situación} & \textbf{Respuesta del firmware}\\
\midrule
MPU6050 no responde al arrancar & Se intenta encontrar de nuevo cada 5 s.\\
Lecturas I2C fallidas & Se descartan y se reinicia la detección de caída libre (una lectura mala no puede formar parte de una caída).\\
Lectura con los tres ejes en cero & Se trata como basura (físicamente imposible salvo error de bus).\\
20 fallas seguidas & Se marca el sensor como perdido y se vuelve a buscar.\\
Errores internos del core de Arduino & Se silencian con \texttt{CORE\_DEBUG\_LEVEL=0} para no inundar el monitor.\\
\bottomrule
\end{tabular}
\end{center}

%=====================================================================
\section{Comunicación con el computador}
%=====================================================================

\subsection{Protocolo}

La ESP32 se conecta a la red WiFi como cliente (modo \emph{station}) y envía cada
evento como un \texttt{HTTP POST} a \texttt{http://<IP del PC>:5000/alerta}, con un cuerpo JSON:

\begin{lstlisting}[numbers=none,language={}]
{"dispositivo":"ESP32-S3","evento":"ALERTA","ip":"192.168.5.110",
 "mac":"E0:72:A1:AA:20:84","causa":"RITMO+CAIDA"}
\end{lstlisting}

\begin{center}
\begin{tabular}{@{}ll@{}}
\toprule
\textbf{Campo} & \textbf{Contenido}\\
\midrule
\texttt{evento} & \texttt{CONECTADA} (al arrancar) o \texttt{ALERTA}\\
\texttt{causa} & Solo en alertas: \texttt{BOTON}, \texttt{RITMO}, \texttt{CAIDA} o combinaciones con \texttt{+}\\
\texttt{ip} & Dirección IP de la placa en la red local\\
\texttt{mac} & Identificador de hardware único de la placa\\
\bottomrule
\end{tabular}
\end{center}

El servidor responde \texttt{200 OK}; ese código es la confirmación de entrega que usa
la placa para decidir si reintenta.

\subsection{Cola de envíos}

El envío no bloquea la vigilancia. Al terminar un conteo, la alerta se agrega a una cola
(máximo 5 elementos) y el sistema vuelve a ESTABLE de inmediato. En cada iteración del
ciclo principal se intenta enviar el primer elemento:

\begin{center}
\begin{tikzpicture}[font=\small,
  caja/.style={draw,minimum width=2.2cm,minimum height=0.8cm,align=center,line width=0.8pt},
  flecha/.style={-{Stealth[length=2.5mm]},thick}]
  \node[caja,draw=cPulso,fill=cPulso!12] (a1) at (0,0) {\scriptsize RITMO+CAIDA};
  \node[caja,draw=cCaida,fill=cCaida!12] (a2) at (2.3,0) {\scriptsize CAIDA};
  \node[caja,draw=gray,fill=gray!8] (a3) at (4.6,0) {\scriptsize (vacío)};
  \node[above=2pt of a2,font=\scriptsize] {cola};
  \node[draw=cPC,fill=cPC!10,rounded corners,minimum width=2.6cm,minimum height=1cm,align=center] (pc) at (-4,0) {Computador};
  \draw[flecha,cPC] (a1.west) -- node[above,font=\scriptsize]{POST} (pc.east);
  \node[align=left,font=\scriptsize,anchor=north west] at (-5.3,-0.9) {
    \textbf{200 OK} $\rightarrow$ se saca de la cola\\
    \textbf{falla} $\rightarrow$ reintenta cada 3 s; tras 5 intentos se descarta y se informa en el monitor};
\end{tikzpicture}
\end{center}

Cada intento tiene un límite de 2 s (\texttt{setConnectTimeout}, \texttt{setTimeout}),
de modo que un computador caído no detiene la lectura de los sensores por mucho tiempo.

\subsection{Servidor en el computador}

\texttt{servidor\_pc.py} usa solo la biblioteca estándar de Python:
\begin{itemize}[leftmargin=*]
  \item \texttt{ThreadingHTTPServer}: atiende cada petición en un hilo propio, así una
        conexión lenta no bloquea a las demás.
  \item \texttt{timeout = 5}: libera conexiones que quedaron abiertas a medias (por
        ejemplo, si el WiFi de la placa se cortó durante un envío).
  \item Responde con \texttt{Content-Length} y \texttt{Connection: close} explícitos,
        para que la placa sepa exactamente cuándo terminó la respuesta.
  \item Muestra en consola una línea para \texttt{CONECTADA} y un recuadro destacado
        (con pitido) para cada \texttt{ALERTA}.
\end{itemize}

%=====================================================================
\section{Estructura del firmware}
%=====================================================================

\subsection{Organización del repositorio}

\begin{lstlisting}[numbers=none,language={}]
alerta-esp32s3/
|-- platformio.ini          configuración de compilación (placa, puerto, flags)
|-- src/main.cpp            firmware
|-- include/config.h        pines, tiempos y umbrales (todo lo ajustable)
|-- include/secrets.h       WiFi y URL del PC (no se sube al repositorio)
|-- pc/servidor_pc.py       servidor que recibe las alertas
`-- docs/                   documentación
\end{lstlisting}

\subsection{Modelo de ejecución}

El framework de Arduino ejecuta \texttt{setup()} una vez y luego llama a \texttt{loop()}
indefinidamente. No hay hilos ni interrupciones: todo es \textbf{cooperativo y no
bloqueante}. Ninguna función espera con \texttt{delay()} durante la operación normal;
los tiempos se controlan comparando con \texttt{millis()} (milisegundos desde el
arranque).

\begin{analogia}[millis()]
Es un \emph{event loop} de un solo hilo, como el de Node.js: cada tarea revisa
``¿ya pasó mi tiempo?'' y devuelve el control rápido. Un \texttt{delay(30000)} sería
equivalente a un \texttt{sleep} síncrono que congela todo el sistema.
\end{analogia}

\subsection{Una iteración de \texttt{loop()}}

\begin{center}
\begin{tikzpicture}[font=\small,node distance=0.35cm,
  f/.style={draw,rounded corners=3pt,minimum width=11cm,align=left,inner sep=5pt,text width=10.6cm,line width=0.7pt},
  flecha/.style={-{Stealth[length=2mm]},thick}]
  \node[f,draw=cBoton,fill=cBoton!8] (1) {\texttt{botonPresionado()} — lee el botón con debounce};
  \node[f,draw=cPulso,fill=cPulso!8,below=of 1] (2) {\texttt{actualizarPulso()} — muestra el ADC cada 10 ms y recalcula BPM};
  \node[f,draw=cCaida,fill=cCaida!8,below=of 2] (3) {\texttt{detectarCaida()} — lee el MPU6050 cada 10 ms y busca la firma de caída};
  \node[f,draw=gray,fill=gray!6,below=of 3] (4) {Reintento de búsqueda del MPU6050 (cada 5 s) y reporte al monitor (cada 2 s)};
  \node[f,draw=cEsp,fill=cEsp!8,below=of 4] (5) {\texttt{switch (estado)} — transiciones de la máquina de estados};
  \node[f,draw=cLed,fill=cLed!12,below=of 5] (6) {\texttt{actualizarLed()} — titila, fijo o apagado según el estado};
  \node[f,draw=cPC,fill=cPC!8,below=of 6] (7) {\texttt{atenderEnvios()} — un intento de envío del primer elemento de la cola};
  \foreach \a/\b in {1/2,2/3,3/4,4/5,5/6,6/7} \draw[flecha] (\a) -- (\b);
  \draw[flecha] (7.east) -- ++(0.6,0) |- node[pos=0.25,right,font=\scriptsize]{siguiente iteración} (1.east);
\end{tikzpicture}
\end{center}

\subsection{Mapa de funciones de \texttt{main.cpp}}

\begin{center}
\begin{tabular}{@{}p{4.4cm} p{11.2cm}@{}}
\toprule
\textbf{Función} & \textbf{Responsabilidad}\\
\midrule
\texttt{actualizarLed()} & Traduce el estado a un patrón del LED.\\
\texttt{textoCausas(m)} & Convierte el conjunto de \emph{flags} (\texttt{C\_BOTON=1}, \texttt{C\_RITMO=2}, \texttt{C\_CAIDA=4}) en texto: \texttt{"RITMO+CAIDA"}.\\
\texttt{botonPresionado()} & Devuelve \texttt{true} solo en el flanco de bajada (el instante de presionar), ya filtrado.\\
\texttt{actualizarPulso()} & Muestreo del ADC, base y amplitud, detección de latidos, cálculo de BPM y filtros.\\
\texttt{ritmoAlto()} & Decide si el pulso dispara, con confirmación temporal e histéresis.\\
\texttt{iniciarMPU()} & Despierta el sensor (registro \texttt{0x6B}) y fija el rango $\pm 8$ g (registro \texttt{0x1C}).\\
\texttt{leerAceleracionG()} & Lee 6 bytes desde el registro \texttt{0x3B} y devuelve $|a|$ en g, o $-1$ si la lectura falla.\\
\texttt{detectarCaida()} & Detección de caída libre + impacto, y tolerancia a fallos del bus.\\
\texttt{iniciarConteo()}\newline\texttt{agregarCausa()} & Entrada a PENDIENTE y suma de causas durante el conteo.\\
\texttt{conectarWiFi()} & Reconexión con límite de 15 s.\\
\texttt{enviarEvento()} & Arma el JSON y hace el POST; devuelve \texttt{true} solo con 200.\\
\texttt{encolarAlerta()}\newline\texttt{atenderEnvios()} & Cola de envíos con reintentos acotados.\\
\texttt{conectarWiFiConLed()} & Conexión inicial con indicación por LED.\\
\texttt{setup()} & Configuración de pines, sensor, WiFi y aviso \texttt{CONECTADA}.\\
\texttt{loop()} & Ciclo principal descrito arriba.\\
\bottomrule
\end{tabular}
\end{center}

\subsection{Fragmentos clave}

\noindent\textbf{Suma de causas con \emph{bit flags}.} Cada causa es un bit; un conjunto de
causas es un OR de bits. Detectar varias a la vez es natural:\par

\begin{lstlisting}[language=C++,numbers=none]
const uint8_t C_BOTON = 1, C_RITMO = 2, C_CAIDA = 4;   // 001, 010, 100
uint8_t nuevas = 0;
if (presion)     nuevas |= C_BOTON;
if (caida)       nuevas |= C_CAIDA;
if (ritmoAlto()) nuevas |= C_RITMO;
if (nuevas) iniciarConteo(nuevas);                     // p. ej. 110 = RITMO+CAIDA
\end{lstlisting}

\noindent\textbf{Fin del conteo: encolar y volver a vigilar.}\par

\begin{lstlisting}[language=C++,numbers=none]
if (estado == PENDIENTE && millis() - tInicioAlerta >= VENTANA_MS) {
  encolarAlerta(causas);                         // el envío ocurre en atenderEnvios()
  tLedApagadoHasta = millis() + LED_APAGADO_MS;  // LED apagado 1 s
  estado = ESTABLE;                              // sigue vigilando de inmediato
}
\end{lstlisting}

\noindent\textbf{Reintentos acotados.}\par

\begin{lstlisting}[language=C++,numbers=none]
void atenderEnvios() {
  if (nCola == 0) return;
  if (intentos > 0 && millis() - tUltimoIntento < REINTENTO_MS) return;
  tUltimoIntento = millis();
  intentos++;
  bool ok = enviarEvento("ALERTA", textoCausas(cola[0]).c_str());
  if (ok || intentos >= MAX_INTENTOS) {   // entregado, o se rinde tras 5 intentos
    /* sacar cola[0] de la cola */
    intentos = 0;
  }
}
\end{lstlisting}

%=====================================================================
\section{Parámetros configurables}
%=====================================================================

Todo lo ajustable está en \texttt{include/config.h}. Un cambio requiere volver a
compilar y cargar el firmware.

\begin{center}
\begin{tabular}{@{}l c p{8.3cm}@{}}
\toprule
\textbf{Parámetro} & \textbf{Valor} & \textbf{Efecto}\\
\midrule
\multicolumn{3}{@{}l}{\textit{Tiempos}}\\
\texttt{VENTANA\_MS} & 30000 & Duración del conteo para cancelar\\
\texttt{DEBOUNCE\_MS} & 50 & Estabilidad exigida al botón\\
\texttt{LED\_TITILEO\_MS} & 500 & Período de medio ciclo del titileo\\
\texttt{LED\_APAGADO\_MS} & 1000 & LED apagado tras enviar una alerta\\
\texttt{REINTENTO\_MS} & 3000 & Espera entre reintentos de envío\\
\texttt{MAX\_INTENTOS} & 5 & Intentos por alerta antes de descartarla\\
\texttt{MUESTREO\_MS} & 10 & Período de muestreo de los sensores (100 Hz)\\
\multicolumn{3}{@{}l}{\textit{Pulso}}\\
\texttt{UMBRAL\_BPM} & 110 & Pulso que dispara la alerta\\
\texttt{PULSO\_AMPLITUD\_MIN} & 25 & Amplitud mínima para considerar que hay señal\\
(filtros) & — & Ver la tabla de filtros anti-ruido en la sección de pulso\\
\multicolumn{3}{@{}l}{\textit{Caídas}}\\
\texttt{CAIDA\_LIBRE\_G} & 0.6 & Debajo de este valor se considera caída libre\\
\texttt{IMPACTO\_G} & 2.0 & Encima de este valor se considera impacto. Subirlo si hay falsos positivos; bajarlo si no detecta caídas\\
\texttt{VENTANA\_IMPACTO\_MS} & 1000 & Tiempo máximo entre caída libre e impacto\\
\multicolumn{3}{@{}l}{\textit{Pines}}\\
\texttt{PIN\_BOTON}, \texttt{PIN\_LED} & 4, 2 & Botón y LED\\
\texttt{PIN\_PULSO} & 5 & Señal del XD-58C (debe ser un pin con ADC1)\\
\texttt{PIN\_SDA}, \texttt{PIN\_SCL} & 8, 9 & Bus I2C del MPU6050\\
\bottomrule
\end{tabular}
\end{center}

\begin{nota}[Por qué el pulso usa un pin de ADC1]
La ESP32 tiene dos convertidores analógicos, ADC1 y ADC2. El ADC2 lo usa internamente el
WiFi, así que mientras el WiFi está activo sus pines no se pueden leer. GPIO5 pertenece
al ADC1 y funciona siempre.
\end{nota}

%=====================================================================
\section{Operación y diagnóstico}
%=====================================================================

\subsection{Puesta en marcha}

\begin{enumerate}[leftmargin=*]
  \item En el computador: \texttt{python pc/servidor\_pc.py}. Muestra la URL en la que escucha.
  \item Poner esa URL en \texttt{PC\_URL} de \texttt{include/secrets.h}, junto con la red y la clave del WiFi (solo redes de 2.4 GHz).
  \item Compilar y cargar con PlatformIO (VS Code $\rightarrow$ \emph{Upload}).
  \item Abrir el monitor serial (\emph{Monitor}) para ver el diagnóstico.
\end{enumerate}

\subsection{Mensajes del monitor serial}

\begin{center}
\begin{tabular}{@{}>{\raggedright\arraybackslash}p{7.3cm} >{\raggedright\arraybackslash}p{7.6cm}@{}}
\toprule
\textbf{Mensaje} & \textbf{Significado}\\
\midrule
\texttt{MPU6050 OK} / \texttt{no encontrado} & Resultado de la detección del acelerómetro al arrancar\\
\texttt{Pulso: 76 BPM | Acel: min 1.00 g, max 1.02 g} & Reporte periódico (cada 2 s) de ambos sensores\\
\texttt{Pulso: sin lectura} & Señal plana: no hay dedo o el sensor no está conectado\\
\texttt{Caída detectada (golpe de 3.1 g)} & Se detectó la firma de caída\\
\texttt{Estado: ALERTA por CAIDA ...} & Comenzó un conteo de 30 s\\
\texttt{Se suma otra causa: ... RITMO+CAIDA} & Otra causa ocurrió durante el conteo\\
\texttt{Cancelado -> Estado: ESTABLE} & El botón canceló el conteo\\
\texttt{30 s cumplidos -> enviando alerta ...} & La alerta entró en la cola de envío\\
\texttt{Evento ALERTA enviado, respuesta HTTP: 200} & Entregada\\
\texttt{respuesta HTTP: -1} & No se pudo conectar con el servidor (se reintenta)\\
\texttt{No se pudo enviar ... después de 5 intentos} & Alerta descartada\\
\bottomrule
\end{tabular}
\end{center}

\subsection{Problemas frecuentes}

\begin{center}
\begin{tabular}{@{}>{\raggedright\arraybackslash}p{5.2cm} >{\raggedright\arraybackslash}p{9.7cm}@{}}
\toprule
\textbf{Síntoma} & \textbf{Causa probable y solución}\\
\midrule
\texttt{MPU6050 no encontrado} o caídas que no se detectan & Cables SDA/SCL flojos o intercambiados. Asegurar la conexión (idealmente soldar).\\
Pulso con valores que saltan sin sentido & Sensor mal conectado o sin alimentación (su LED verde apagado). Revisar + y $-$; activar los filtros anti-ruido cuando el hardware esté firme.\\
\texttt{respuesta HTTP: -1} constante & El servidor no está corriendo, la IP del PC cambió, o el firewall de Windows bloquea el puerto 5000.\\
El monitor no muestra nada & El puerto USB usado no coincide con \texttt{ARDUINO\_USB\_CDC\_ON\_BOOT} en \texttt{platformio.ini} (0 = puerto UART, 1 = USB nativo).\\
LED parpadea rápido al arrancar & No pudo conectarse al WiFi: revisar red y clave en \texttt{secrets.h}.\\
\bottomrule
\end{tabular}
\end{center}

%=====================================================================
\section{Glosario}
%=====================================================================

\begin{description}[leftmargin=*,style=nextline]
  \item[ADC] Conversor analógico-digital: mide un voltaje y lo entrega como número entero.
  \item[BPM] Pulsaciones por minuto.
  \item[Debounce] Filtrado de los rebotes mecánicos de un botón.
  \item[Firmware] Programa grabado en la memoria del microcontrolador.
  \item[g] Unidad de aceleración igual a la gravedad terrestre (9.8 m/s$^2$).
  \item[GND] Tierra, referencia de 0 V común a todo el circuito.
  \item[GPIO] Pin de propósito general del microcontrolador.
  \item[I2C] Bus serial de dos cables (SDA datos, SCL reloj) para comunicar chips.
  \item[IBI] Intervalo entre latidos, en milisegundos.
  \item[Microcontrolador] Computador completo en un solo chip (procesador, memoria, periféricos).
  \item[Pull-up] Resistencia que mantiene un pin en \texttt{HIGH} cuando nada lo conecta a GND.
  \item[UART] Puerto serial; aquí, el canal por el que la placa envía texto al monitor.
  \item[3V3] Alimentación de 3.3 V que la ESP32 entrega a los sensores.
\end{description}

\appendix
%=====================================================================
\section{Código fuente completo}
%=====================================================================

\subsection{\texttt{src/main.cpp}}
\begin{lstlisting}[language=C++]

#include <Arduino.h>
#include <WiFi.h>
#include <HTTPClient.h>
#include <Wire.h>

#include "config.h"
#include "secrets.h"

// ESTABLE = vigilando, PENDIENTE = contando 30 s para cancelar.
// Los sensores se leen SIEMPRE, en cualquier estado, también después de enviar una alerta.
enum Estado { ESTABLE, PENDIENTE };
Estado estado = ESTABLE;

unsigned long tInicioAlerta = 0;
unsigned long tLedApagadoHasta = 0;  // después de enviar una alerta, el LED se apaga un momento

// LED: titila mientras vigila, fijo mientras cuenta los 30 s,
// y se apaga un momento cuando se envía la alerta.
void actualizarLed() {
  if (estado == PENDIENTE) {
    digitalWrite(PIN_LED, HIGH);                        // contando: fijo
  } else if (millis() < tLedApagadoHasta) {
    digitalWrite(PIN_LED, LOW);                         // alerta enviada: apagado un momento
  } else {
    digitalWrite(PIN_LED, (millis() / LED_TITILEO_MS) % 2 ? HIGH : LOW);  // vigilando: titila
  }
}
// Causas de la alerta: se pueden sumar (por ejemplo RITMO + CAIDA)
const uint8_t C_BOTON = 1, C_RITMO = 2, C_CAIDA = 4;
uint8_t causas = 0;

String textoCausas(uint8_t m) {
  String t = "";
  if (m & C_BOTON) t += "BOTON";
  if (m & C_RITMO) t += (t.length() ? "+RITMO" : "RITMO");
  if (m & C_CAIDA) t += (t.length() ? "+CAIDA" : "CAIDA");
  return t;
}

// Alertas listas para enviar (fila). El envío se hace por detrás, sin dejar de vigilar.
const int MAX_COLA = 5;
uint8_t cola[MAX_COLA];
int nCola = 0;
int intentos = 0;                // intentos de envío de la primera alerta de la fila
unsigned long tUltimoIntento = 0;

// Debounce
int lecturaAnterior = HIGH;
int estadoBoton = HIGH;
unsigned long tCambio = 0;

bool botonPresionado() {
  int lectura = digitalRead(PIN_BOTON);
  if (lectura != lecturaAnterior) tCambio = millis();
  lecturaAnterior = lectura;

  if (millis() - tCambio > DEBOUNCE_MS && lectura != estadoBoton) {
    estadoBoton = lectura;
    if (estadoBoton == LOW) return true;  
  }
  return false;
}

// =====================================================================
//  Sensor de pulso XD-58C (salida analógica)
// =====================================================================
int bpm = 0;                     // pulsaciones por minuto (0 = sin lectura)
float lineaBase = 2048;          // valor promedio de la señal
float amplitud = 0;              // tamaño típico de la onda de pulso
bool sobreUmbral = false;        // la onda está en su parte alta
unsigned long tUltimoLatido = 0;
unsigned long tMuestraPulso = 0;
unsigned long ibis[5];           // tiempos entre latidos (ms) de los últimos 5 latidos
int nIbis = 0, iIbi = 0;
bool ritmoArmado = true;         // false = ya avisó; espera a que el pulso baje
const char* estadoPulso = "sin lectura (¿dedo bien puesto?)";
unsigned long tRitmoAltoDesde = 0;  // desde cuándo el pulso está alto (0 = no está alto)

void actualizarPulso() {
  if (millis() - tMuestraPulso < MUESTREO_MS) return;
  tMuestraPulso = millis();
  unsigned long ahora = tMuestraPulso;

  int v = analogRead(PIN_PULSO);
  lineaBase += (v - lineaBase) * 0.005f;                // promedio lento
  amplitud  += (fabsf(v - lineaBase) - amplitud) * 0.01f;

  // Sensor apagado o desconectado: la señal se queda pegada abajo o arriba
  if (lineaBase < PULSO_BASE_MIN || lineaBase > PULSO_BASE_MAX) {
    bpm = 0; nIbis = 0; sobreUmbral = false;
    estadoPulso = "sensor apagado o desconectado (¿prende su luz verde?)";
    return;
  }
  if (amplitud < PULSO_AMPLITUD_MIN) {                  // señal plana: sin dedo
    bpm = 0; nIbis = 0; sobreUmbral = false;
    estadoPulso = "sin lectura (¿dedo bien puesto?)";
    return;
  }

  // Latido = la señal sube por encima de la mitad de la onda
  if (!sobreUmbral && v > lineaBase + amplitud * 0.5f && ahora - tUltimoLatido > 300) {
    sobreUmbral = true;
    unsigned long ibi = ahora - tUltimoLatido;
    tUltimoLatido = ahora;
    if (ibi >= IBI_MIN_MS && ibi <= IBI_MAX_MS) {       // entre 40 y 180 BPM
      ibis[iIbi] = ibi;
      iIbi = (iIbi + 1) % 5;
      if (nIbis < 5) nIbis++;
    } else {
      nIbis = 0;                                        // latido raro: empezar a contar de nuevo
    }
    if (nIbis == 5) {
      unsigned long suma = 0, minIbi = ibis[0], maxIbi = ibis[0];
      for (int i = 0; i < 5; i++) {
        suma += ibis[i];
        if (ibis[i] < minIbi) minIbi = ibis[i];
        if (ibis[i] > maxIbi) maxIbi = ibis[i];
      }
      // Un corazón late parejo; el ruido da tiempos muy distintos entre "latidos"
      if (maxIbi > minIbi * PULSO_REGULARIDAD) {
        bpm = 0;
        estadoPulso = "señal irregular (ruido), se ignora";
      } else {
        bpm = 60000UL * 5 / suma;                       // promedio de 5 latidos
      }
    }
  } else if (sobreUmbral && v < lineaBase) {
    sobreUmbral = false;
  }

  if (ahora - tUltimoLatido > 2000) {                   // 2 s sin latidos
    bpm = 0; nIbis = 0;
    estadoPulso = "sin lectura (¿dedo bien puesto?)";
  }
}

// true cuando el pulso lleva RITMO_CONFIRMA_MS seguidos en UMBRAL_BPM o más.
// Después de avisar, vuelve a avisar solo cuando el pulso baje de UMBRAL_BPM - 5.
bool ritmoAlto() {
  if (bpm > 0 && bpm < UMBRAL_BPM - 5) ritmoArmado = true;
  if (bpm < UMBRAL_BPM) { tRitmoAltoDesde = 0; return false; }
  if (tRitmoAltoDesde == 0) tRitmoAltoDesde = millis();
  return ritmoArmado && millis() - tRitmoAltoDesde >= RITMO_CONFIRMA_MS;
}

// =====================================================================
//  Acelerómetro MPU6050 (I2C) - detección de caídas
// =====================================================================
const uint8_t MPU_ADDR = 0x68;
bool mpuOk = false;
bool enCaidaLibre = false;
unsigned long tCaidaLibre = 0;
unsigned long tMuestraMPU = 0;
float gMin = 99, gMax = 0;       // aceleración mínima y máxima vista (para revisar en el monitor)
int fallosMPU = 0;               // lecturas fallidas seguidas (cables flojos)

bool iniciarMPU() {
  Wire.begin(PIN_SDA, PIN_SCL);
  Wire.beginTransmission(MPU_ADDR);
  Wire.write(0x6B); Wire.write(0x00);                   // despertar el sensor
  if (Wire.endTransmission() != 0) return false;
  Wire.beginTransmission(MPU_ADDR);
  Wire.write(0x1C); Wire.write(0x10);                   // rango ±8 g
  Wire.endTransmission();
  return true;
}

// Devuelve la aceleración total en g (1 g = quieto). -1 si falla la lectura.
float leerAceleracionG() {
  Wire.beginTransmission(MPU_ADDR);
  Wire.write(0x3B);
  if (Wire.endTransmission(false) != 0) return -1;
  if (Wire.requestFrom(MPU_ADDR, (uint8_t)6) != 6) return -1;
  uint8_t b[6];
  for (int i = 0; i < 6; i++) b[i] = Wire.read();       // leer en orden
  int16_t ax = (b[0] << 8) | b[1];
  int16_t ay = (b[2] << 8) | b[3];
  int16_t az = (b[4] << 8) | b[5];
  if (ax == 0 && ay == 0 && az == 0) return -1;         // todo en cero = lectura basura
  float x = ax / 4096.0f, y = ay / 4096.0f, z = az / 4096.0f;  // ±8 g -> 4096 por g
  return sqrtf(x * x + y * y + z * z);
}

// Caída = caída libre (casi 0 g) seguida de un golpe fuerte en menos de 1 s.
bool detectarCaida() {
  if (!mpuOk || millis() - tMuestraMPU < MUESTREO_MS) return false;
  tMuestraMPU = millis();
  float g = leerAceleracionG();
  if (g < 0) {
    enCaidaLibre = false;                               // con una lectura mala no se confía
    if (++fallosMPU >= 20) {
      mpuOk = false;
      fallosMPU = 0;
      Serial.println("MPU6050 perdió la conexión (revisa SDA/SCL): lo vuelvo a buscar");
    }
    return false;
  }
  fallosMPU = 0;
  if (g < gMin) gMin = g;
  if (g > gMax) gMax = g;

  if (g < CAIDA_LIBRE_G) {
    if (!enCaidaLibre) { enCaidaLibre = true; tCaidaLibre = millis(); }
  } else if (enCaidaLibre && g > IMPACTO_G) {
    enCaidaLibre = false;
    Serial.printf("Caída detectada (golpe de %.1f g)\n", g);
    return true;
  }
  if (enCaidaLibre && millis() - tCaidaLibre > VENTANA_IMPACTO_MS) enCaidaLibre = false;
  return false;
}

// Pasa a PENDIENTE: LED encendido y empieza el conteo de 30 s
void iniciarConteo(uint8_t nuevas) {
  estado = PENDIENTE;
  causas = nuevas;
  tInicioAlerta = millis();
  digitalWrite(PIN_LED, HIGH);
  Serial.printf("Estado: ALERTA por %s (30 s para cancelar con el botón)\n", textoCausas(causas).c_str());
}

// Durante el conteo se puede sumar otra causa (por ejemplo, una caída mientras el pulso está alto)
void agregarCausa(uint8_t c) {
  if (causas & c) return;
  causas |= c;
  Serial.printf("Se suma otra causa: ahora la alerta es por %s\n", textoCausas(causas).c_str());
}

void conectarWiFi() {
  if (WiFi.status() == WL_CONNECTED) return;
  WiFi.begin(WIFI_SSID, WIFI_PASS);
  Serial.print("Conectando WiFi");
  unsigned long t0 = millis();
  while (WiFi.status() != WL_CONNECTED && millis() - t0 < 15000) {
    delay(250);
    Serial.print(".");
  }
  Serial.println(WiFi.status() == WL_CONNECTED ? " OK" : " FALLÓ");
  if (WiFi.status() == WL_CONNECTED) {
    Serial.print("IP: ");
    Serial.println(WiFi.localIP());
  }
}

bool enviarEvento(const char* evento, const char* causa = "") {
  conectarWiFi();
  if (WiFi.status() != WL_CONNECTED) return false;

  HTTPClient http;
  http.setConnectTimeout(2000);   // máximo 2 s esperando al PC, para no dejar de vigilar
  http.setTimeout(2000);
  http.begin(PC_URL);
  http.addHeader("Content-Type", "application/json");
  String body = String("{\"dispositivo\":\"ESP32-S3\",\"evento\":\"") + evento +
                "\",\"ip\":\"" + WiFi.localIP().toString() +
                "\",\"mac\":\"" + WiFi.macAddress() + "\"";
  if (causa[0] != '\0') body += String(",\"causa\":\"") + causa + "\"";
  body += "}";
  int code = http.POST(body);
  http.end();

  Serial.printf("Evento %s enviado, respuesta HTTP: %d\n", evento, code);
  return code == 200;
}

// Pone una alerta en la fila para enviarla
void encolarAlerta(uint8_t m) {
  if (nCola < MAX_COLA) cola[nCola++] = m;
  else Serial.println("Fila de alertas llena: se descarta la más nueva");
}

// Envía la primera alerta de la fila. Si falla, reintenta cada REINTENTO_MS,
// máximo MAX_INTENTOS veces, para no quedarse mandando mensajes sin parar.
void atenderEnvios() {
  if (nCola == 0) return;
  if (intentos > 0 && millis() - tUltimoIntento < REINTENTO_MS) return;
  tUltimoIntento = millis();
  intentos++;
  String c = textoCausas(cola[0]);
  bool ok = enviarEvento("ALERTA", c.c_str());
  if (ok || intentos >= MAX_INTENTOS) {
    if (!ok) Serial.printf("No se pudo enviar la alerta %s después de %d intentos\n", c.c_str(), intentos);
    for (int i = 1; i < nCola; i++) cola[i - 1] = cola[i];   // sacar de la fila
    nCola--;
    intentos = 0;
  }
}

// Parpadea el LED mientras intenta conectar al WiFi (máx. 15 s).
// Al terminar: conectado = LED fijo 2 s; sin WiFi = 10 parpadeos rápidos ( esto solo lo puse porque no sabia si me estaba leyendo:)).
void conectarWiFiConLed() {
  WiFi.begin(WIFI_SSID, WIFI_PASS);
  Serial.print("Conectando WiFi");
  unsigned long t0 = millis();
  while (WiFi.status() != WL_CONNECTED && millis() - t0 < 15000) {
    digitalWrite(PIN_LED, !digitalRead(PIN_LED));
    delay(250);
    Serial.print(".");
  }
  if (WiFi.status() == WL_CONNECTED) {
    Serial.printf(" OK\nIP: %s\n", WiFi.localIP().toString().c_str());
    digitalWrite(PIN_LED, HIGH);
    delay(2000);
  } else {
    Serial.println(" FALLÓ");
    for (int i = 0; i < 20; i++) {
      digitalWrite(PIN_LED, !digitalRead(PIN_LED));
      delay(80);
    }
  }
  digitalWrite(PIN_LED, LOW);
}

void setup() {
  Serial.begin(115200);
  unsigned long t0 = millis();
  while (!Serial && millis() - t0 < 3000) delay(10);  // espera al monitor USB

  pinMode(PIN_BOTON, INPUT_PULLUP);
  pinMode(PIN_LED, OUTPUT);
  digitalWrite(PIN_LED, LOW);
  analogReadResolution(12);                 // pulso: valores de 0 a 4095
  mpuOk = iniciarMPU();
  Serial.println(mpuOk ? "MPU6050 OK" : "MPU6050 no encontrado: revisa SDA/SCL (caídas desactivadas)");
  WiFi.mode(WIFI_STA);
  WiFi.setAutoReconnect(true);
  conectarWiFiConLed();
  enviarEvento("CONECTADA");  // aviso al PC de que la placa está en línea
  Serial.println("Estado: ESTABLE");
}

void loop() {
  bool presion = botonPresionado();
  actualizarPulso();
  bool caida = detectarCaida();

  // Si el MPU6050 no respondió (cables flojos), lo vuelve a buscar cada 5 s
  static unsigned long tBuscarMPU = 0;
  if (!mpuOk && millis() - tBuscarMPU >= 5000) {
    tBuscarMPU = millis();
    mpuOk = iniciarMPU();
    if (mpuOk) Serial.println("MPU6050 encontrado");
  }

  // Cada 2 s muestra el pulso y la aceleración en el monitor (para revisar los sensores)
  static unsigned long tImpresion = 0;
  if (millis() - tImpresion >= 2000) {
    tImpresion = millis();
    if (bpm > 0) Serial.printf("Pulso: %d BPM", bpm);
    else Serial.printf("Pulso: %s", estadoPulso);
    if (mpuOk) Serial.printf("  |  Acel: min %.2f g, max %.2f g\n", gMin, gMax);
    else Serial.printf("  |  Acel: MPU6050 no encontrado (revisa SDA/SCL)\n");
    gMin = 99; gMax = 0;
  }

  switch (estado) {
    case ESTABLE: {
      // Se revisan las tres causas a la vez: si pasan juntas, la alerta las lleva todas
      uint8_t nuevas = 0;
      if (presion) nuevas |= C_BOTON;
      if (caida)   nuevas |= C_CAIDA;
      if (ritmoAlto()) {
        ritmoArmado = false;          // no volver a avisar hasta que el pulso baje
        nuevas |= C_RITMO;
      }
      if (nuevas) iniciarConteo(nuevas);
      break;
    }

    case PENDIENTE:
      if (presion) {
        estado = ESTABLE;
        Serial.println("Cancelado -> Estado: ESTABLE");
      } else {
        // Si durante el conteo pasa otra cosa, se suma a la alerta
        if (caida) agregarCausa(C_CAIDA);
        if (ritmoAlto()) { ritmoArmado = false; agregarCausa(C_RITMO); }
      }
      if (estado == PENDIENTE && millis() - tInicioAlerta >= VENTANA_MS) {
        Serial.printf("30 s cumplidos -> enviando alerta por %s (y sigo vigilando)\n", textoCausas(causas).c_str());
        encolarAlerta(causas);
        tLedApagadoHasta = millis() + LED_APAGADO_MS;  // el LED se apaga un momento
        estado = ESTABLE;             // vuelve a vigilar (LED titilando): otra caída o pulso alto manda otra alerta
      }
      break;
  }

  actualizarLed();
  atenderEnvios();
}
\end{lstlisting}

\subsection{\texttt{include/config.h}}
\begin{lstlisting}[language=C++]
#pragma once

// ---------- Pines ----------
constexpr int PIN_BOTON = 4;   // botón entre GPIO4 y GND (pull-up interno)
constexpr int PIN_LED   = 2;   // LED + resistencia (220-330 Ω) a GND

// ---------- Tiempos ----------
constexpr unsigned long VENTANA_MS   = 30000;  // 30 s para cancelar la alerta
constexpr unsigned long DEBOUNCE_MS  = 50;
constexpr unsigned long LED_TITILEO_MS = 500;   // vigilando: el LED cambia cada 0.5 s
constexpr unsigned long LED_APAGADO_MS = 1000;  // al enviar una alerta: apagado 1 s
constexpr unsigned long REINTENTO_MS = 3000;   // reintento si falla el envío
constexpr int MAX_INTENTOS = 5;                 // si falla 5 veces, deja de reintentar esa alerta

// ---------- Sensores ----------
constexpr int PIN_PULSO = 5;   // XD-58C: pata S (señal) -> GPIO5 (entrada analógica)
constexpr int PIN_SDA   = 8;   // MPU6050: SDA -> GPIO8
constexpr int PIN_SCL   = 9;   // MPU6050: SCL -> GPIO9

constexpr unsigned long MUESTREO_MS = 10;  // leer los sensores cada 10 ms

// Ritmo cardíaco
constexpr int UMBRAL_BPM         = 110;    // desde aquí empieza el conteo de 30 s
constexpr int PULSO_AMPLITUD_MIN = 25;     // por debajo = no hay dedo / señal plana

// Filtros anti-ruido: APAGADOS (lee como al inicio, aunque sea con ruido).
// Para prenderlos: BASE 300/3800, REGULARIDAD 1.3, IBI 333/1500, CONFIRMA 5000.
constexpr int PULSO_BASE_MIN      = 0;      // 0 = no revisa si el sensor está apagado
constexpr int PULSO_BASE_MAX      = 4096;   // 4096 = no revisa si el sensor está apagado
constexpr float PULSO_REGULARIDAD = 100.0f; // 100 = acepta latidos disparejos
constexpr unsigned long IBI_MIN_MS = 300;   // 200 BPM
constexpr unsigned long IBI_MAX_MS = 2000;  // 30 BPM
constexpr unsigned long RITMO_CONFIRMA_MS = 0;  // 0 = dispara apenas llega a UMBRAL_BPM

// Caídas
constexpr float CAIDA_LIBRE_G             = 0.6f;  // casi sin peso = está cayendo
constexpr float IMPACTO_G                 = 2.0f;  // golpe fuerte = llegó al suelo
constexpr unsigned long VENTANA_IMPACTO_MS = 1000; // el golpe debe llegar en menos de 1 s
\end{lstlisting}

\subsection{\texttt{pc/servidor\_pc.py}}
\begin{lstlisting}[language=Python]


import json
import socket
from datetime import datetime
from http.server import BaseHTTPRequestHandler, ThreadingHTTPServer

PUERTO = 5000


class Handler(BaseHTTPRequestHandler):
    timeout = 5  # si una conexión se queda colgada, se suelta a los 5 s

    def do_POST(self):
        if self.path != "/alerta":
            self.send_response(404)
            self.end_headers()
            return
        largo = int(self.headers.get("Content-Length", 0))
        datos = self.rfile.read(largo).decode(errors="ignore")
        try:
            datos = json.loads(datos)
        except json.JSONDecodeError:
            pass
        hora = datetime.now().strftime("%Y-%m-%d %H:%M:%S")
        evento = datos.get("evento") if isinstance(datos, dict) else None
        if evento == "CONECTADA":
            print(f"[{hora}] ESP32 conectada al WiFi desde {self.client_address[0]}")
        else:
            print("\a")  # pitido de consola
            print("=" * 50)
            print(f"  ¡ALERTA RECIBIDA!  {hora}")
            print(f"  Desde: {self.client_address[0]}  Datos: {datos}")
            print("=" * 50)
        self.send_response(200)
        self.send_header("Content-Length", "2")
        self.send_header("Connection", "close")
        self.end_headers()
        self.wfile.write(b"OK")

    def log_message(self, *args):
        pass  # silencia el log por defecto


if __name__ == "__main__":
    s = socket.socket(socket.AF_INET, socket.SOCK_DGRAM)
    try:
        s.connect(("8.8.8.8", 80))
        ip = s.getsockname()[0]
    except OSError:
        ip = "desconocida"
    finally:
        s.close()
    print(f"Escuchando en http://{ip}:{PUERTO}/alerta")
    print("Pon esa IP en PC_URL del sketch del ESP32.")
    ThreadingHTTPServer(("0.0.0.0", PUERTO), Handler).serve_forever()  # atiende varios mensajes a la vez
\end{lstlisting}

\subsection{\texttt{platformio.ini}}
\begin{lstlisting}[language=ini]
[env:esp32-s3-devkitc-1]
platform = espressif32
board = esp32-s3-devkitc-1
framework = arduino
monitor_speed = 115200
; Serial sale por el puerto UART (COM). Para usar el puerto USB nativo del S3, cambiar a 1
build_flags =
    -D ARDUINO_USB_CDC_ON_BOOT=0
    -D CORE_DEBUG_LEVEL=0   ; no llenar el monitor con errores internos (por ejemplo del I2C)
; Evita que el monitor reinicie o bloquee la placa por el USB nativo
monitor_rts = 0
monitor_dtr = 0
\end{lstlisting}



\end{document}
